\section{Introduction}

Computerized interpretation of the 12-lead electrocardiogram (ECG) has long combined signal processing, which can be used to detect
beats and measure intervals, axis and amplitudes, with diagnostic rules applied to those
measurements~\citep{schlapfer2017computer,willems1991diagnostic}. A second classical route keeps the preprocessing and
the measured features but replaces the rules with a statistical or machine-learning classifier trained on labeled
examples~\citep{dechazal2004heartbeats,lyon2018computational}. On the other hand, deep neural networks trained on large
labeled datasets learn the whole mapping from signal to diagnosis, from arrhythmia detection in single-lead
recordings~\citep{hannun2019cardiologist} to multi-label diagnosis of the 12-lead ECG~\citep{ribeiro2020automatic}. PTB-XL, a large public dataset consisting of \NRecordsTotal{} clinical 10-second recordings
from \NPatients{} patients~\citep{Wagner2020,Wagner2022physionet}, is a widely used benchmark for such models, with a
fixed train--test split, label definition and evaluation metric~\citep{Strodthoff2021}. On its diagnostic superclass task, which assigns each
recording to one or more of five classes (normal ECG -- NORM; myocardial infarction -- MI; ST/T change -- STTC; conduction
disturbance -- CD; hypertrophy -- HYP), convolutional and recurrent networks trained on the eight training folds by the
benchmark's authors (the original benchmark models) reached a
macro-averaged area under the receiver operating characteristic curve (AUROC) of up to \BestSingleBenchmarkAUROC, and
their ensemble \BestPublishedAUROC~\citep{Strodthoff2021}. Later supervised and pretrained models report
\MinLaterAUROC--\MaxReportedAUROC~\citep{shi2023multires,zhang2023spacetime,behrouz2024chimera,karami2025msssm,kim2024ecgjepa,lunelli2025benchecg,coppola2024hubertecg,almasud2026benchmarking,almasud2026pretraining,deng2026recursiveecg};
the highest value is from a supervised classifier whose design was revised by a large language model (LLM) agent and evaluated on a
different test set~\citep{deng2026recursiveecg}. Each of these models needs thousands of labeled recordings and returns class scores
without clear interpretation.

LLMs can apply medical diagnostic criteria written in plain language and state reasons for their
output, but reading the ECG signal has been their weak point. Namely, ECG--language models, which pretrain an ECG encoder
on hundreds of thousands of ECGs with their reports or labels extracted from them, reach a macro AUROC of
\MinZeroShotAUROC--\BestZeroShotAUROC{} on the PTB-XL superclass test fold without PTB-XL
labels~\citep{liu2024merl,wang2025melp,pham2025dbeta,cai2025supreme}. Moreover, general-purpose multimodal LLMs that
were shown ECG images scored close to chance, which for the AUROC is 0.5 for every class whatever its prevalence
(\MinGeneralLLMAUROC--\BestGeneralLLMAUROC{} on a \NOtherTestSet-record
PTB-XL hold-out set)~\citep{liu2024teach,lan2025gem}. On the same hold-out set, multimodal LLMs instruction-tuned on ECG
data, that is, further trained on ECG images paired with instructions and expert answers, reached
\MinECGLLMAUROC--\BestECGLLMAUROC~\citep{liu2024teach,lan2025gem}. Another study reports that a general-purpose model
labeled almost all abnormal PTB-XL images as normal~\citep{seki2025assessing}. Newer general-purpose models, when given
12-lead ECG images, still scored near chance in telling normal from abnormal recordings of the Shandong Provincial
Hospital database (AUROC 0.51--0.53)~\citep{ahammad2026ecgimage} or reached a balanced accuracy of at most 62\% on
clinical ECGs~\citep{soubh2026performance}. GPT-5.4, Claude Opus 4.7 and Gemini 3.1 Pro recalled fewer than one in ten
critical findings on PTB-XL~\citep{gao2026luminaecg}. Moreover, given numeric series written as text, general-purpose LLMs scored
no better than chance, or barely above it, on ECG rhythm and question-answering
tasks~\citep{liu2025picture,langer2025opentslm}. Better results needed signal processing outside the model, such as
measurements from ECG software given to the model~\citep{yu2023zeroshot} or a task-specific procedure with a worked
example~\citep{xu2024penetrative}.

A different route is to instruct the language model to do the signal processing itself. The motivation for this approach stems
from the fact that LLMs can solve problems by generating executable programs whose execution is delegated to an external
runtime, such as a Python interpreter~\citep{gao2023pal}. This improves numerical reasoning~\citep{chen2023pot} and the
behavior of LLM agents~\citep{wang2024codeact}. Language models can also learn to call external
tools, such as a calculator or a search engine~\citep{schick2023toolformer}; for the ECG, such a tool can be a
measurement library or a pretrained classifier. Commercial LLM interfaces, including the one used
here~\citep{anthropic2026codeexec}, now offer a sandbox in which the model runs Python on files attached to a request.
With a sandbox, a general-purpose LLM could do what ECG software does---filter the raw signal, detect beats, form median
beats and measure them---and then apply written diagnostic criteria to its own measurements, without ECG-specific
training, and its code and printed measurements would leave a record of how each answer was reached. Code-executing LLM
agents have analyzed tabular clinical data~\citep{tayebiarasteh2024automl} and aggregated wearable
data~\citep{merrill2026phia,heydari2025pha}. For physiological waveforms, general-purpose LLMs with code access have
been evaluated on heart-rate estimation from photoplethysmography~\citep{feli2025ppgagent} and, within a broad
time-series benchmark, on single-lead ECG from sleep studies, where they fell well short of specialist
models~\citep{li2026hearts}. Code-writing LLMs have also been used to verify other models' reasoning about 12-lead
ECGs~\citep{xu2026ecgreasoning} and to refine supervised ECG classifiers~\citep{deng2026recursiveecg}, and ECG agents
were used to call fixed, pretrained tools: a self-supervised classifier, NeuroKit2 measurements and an explanation
module~\citep{chung2026ecgagent}. Zero-shot LLM evaluations on PTB-XL have supplied
precomputed features~\citep{yu2023zeroshot,yu2025alfred} or rendered ECG
images~\citep{liu2024teach,seki2025assessing,gao2026luminaecg}.

To our knowledge, no study has evaluated a
general-purpose LLM that writes and executes its own signal-processing code on raw 12-lead ECGs as a zero-shot
diagnostic classifier on the PTB-XL superclass benchmark. Here, we evaluate this approach. We ask three questions. (i) How well does a general-purpose LLM that is instructed to write and run
its own signal-processing code classify raw 12-lead ECGs into the PTB-XL superclasses without ECG-specific training,
compared with published supervised, pretrained and zero-shot models and with a control that uses only the subject's age and sex information?
(ii) What does the model do in the sandbox, and how does it arrive at the measurements it states? (iii) How well do its
superclass probabilities separate records by individual diagnostic statement, and are the measurements it states
consistent with the dataset's annotations? The first question was pre-specified, while the other two were examined post hoc.

The methodological setup can be summarized as follows. A general-purpose LLM received each raw 10-second, 12-lead ECG
recording at 500\,Hz as a file, together with age and sex,
and a prompt that asked it to analyze the ECG in a Python sandbox and then to assign a probability to each PTB-XL
superclass. The model received no labeled examples, no class prevalences and no other external information, and its
parameters were not updated at any stage. The configuration was developed on \NDevRecordsUsed{} of the \NValFold{}
records of the validation fold and fixed in a protocol that was frozen before a single run of this configuration on all
\NTest{} records of the benchmark test fold.

This work makes three contributions: (1) an approach in which a general-purpose LLM agent, instead of calling fixed
tools, writes and executes its own signal-processing code on raw 12-lead ECGs and acts as a zero-shot classifier, and its
evaluation on the PTB-XL superclass benchmark under a protocol frozen before the test run; (2) a description of what the model did in the sandbox, how it obtained
the measurements it stated, and how its probabilities and measurements relate to the dataset's diagnostic statements
and annotations; and (3) a way to audit such an agent: the release of every request, sandbox transcript, printed output and answer, so
that each prediction can be traced to the code and measurements it rests on, together with an offline replay of all of
the model's code runs, which checks that the printed outputs can be reproduced without the API.
